\documentclass[10pt]{article}
\usepackage[letterpaper,margin=0.68in]{geometry}
\usepackage{amsmath,amssymb}
\usepackage{enumitem}
\usepackage{xcolor}
\usepackage{fancyhdr}
\usepackage{microtype}
\usepackage[most]{tcolorbox}
\definecolor{olive}{HTML}{4D5430}
\definecolor{gold}{HTML}{9A7B22}
\definecolor{pale}{HTML}{F5F1DC}
\definecolor{softgreen}{HTML}{EDF0E2}
\pagestyle{fancy}\fancyhf{}\lhead{\textcolor{olive}{MATH\& 107 Answer Key}}\rhead{\textcolor{gold}{Fall 2026}}\cfoot{\thepage}
\setlength{\headheight}{14pt}
\setlength{\parindent}{0pt}\setlength{\parskip}{5pt}
\setlist[enumerate]{leftmargin=*,itemsep=5pt,topsep=4pt}
\newtcolorbox{solbox}[1]{colback=pale,colframe=olive,boxrule=.7pt,arc=2mm,title=\textbf{#1},fonttitle=\color{olive},left=2mm,right=2mm,top=1.5mm,bottom=1.5mm}
\newcommand{\modulekey}[2]{\clearpage{\Large\bfseries\textcolor{olive}{Module #1: #2 — Answer Key}}\par\textcolor{gold}{\rule{\linewidth}{1.2pt}}\par}
\begin{document}
\modulekey{10}{Fractal Geometry and the Limits of Prediction}
\textbf{Prequiz}
\begin{enumerate}
\item
\[
7\to22\to11\to34\to17\to52\to26\to13\to40\to20\to10\to5\to16\to8\to4\to2\to1.
\]
There are \(\boxed{16}\) steps, and the peak value is \(\boxed{52}\).
\item At step \(3\), the number of segments is
\[
4^3=\boxed{64}.
\]
Total length:
\[
\left(\frac43\right)^3=\frac{64}{27}\approx\boxed{2.370}.
\]
The length factor per step is \(\boxed{4/3}\).
\item
\[
\left(\frac34\right)^4=\frac{81}{256}\approx\boxed{0.3164}.
\]
\item
\[
x_1=2.8(0.2)(0.8)=\boxed{0.448}.
\]
To determine \(x_2\), substitute \(x_1=0.448\) into the same recursion with \(r=2.8\).
\item Measured lengths are
\[
10(50)=500,\quad5(130)=650,\quad2.5(340)=850,\quad1.25(900)=1125.
\]
The measured length increases as the ruler gets shorter because the measurement scale becomes finer and captures more small boundary detail.
\end{enumerate}

\textbf{Matching quiz}
\begin{enumerate}
\item
\[
6\to3\to10\to5\to16\to8\to4\to2\to1.
\]
There are \(\boxed8\) steps, and the peak is \(\boxed{16}\).
\item Step \(4\) has \(4^4=\boxed{256}\) segments. Total length is multiplied by
\[
\left(\frac43\right)^4=\frac{256}{81}\approx\boxed{3.1605}
\]
from step \(0\) to step \(4\).
\item
\[
\left(\frac34\right)^3=\frac{27}{64}=\boxed{0.421875}.
\]
\item
\[
x_1=3(0.5)(0.5)=\boxed{0.75},
\]
\[
x_2=3(0.75)(0.25)=\boxed{0.5625}.
\]
\item It demonstrates \textbf{scale-dependent measurement}: a finer ruler follows more small-scale boundary detail, so the numerical measurement can change even though the physical coastline has not changed.
\end{enumerate}
\end{document}
